\chapter*{Conclusion and Perspectives}
\addcontentsline{toc}{chapter}{Conclusion and Perspectives}
\markboth{Conclusion}{Conclusion}
\label{chap:conclusion}

\section*{Summary of Work}

\noindent
This report presented a systematic experimental study of AQ-NBEATS
enhancements for probabilistic short-term electricity demand forecasting
on the EMHIRES MHLV benchmark, covering 35 European countries over a
48-hour forecast horizon. Starting from the published baseline CRPS of
211.22, we designed and evaluated more than twenty configurations across
four complementary categories: exogenous feature and series-embedding
augmentation, curriculum quantile training, novel architectural
contributions, and post-hoc conformal calibration. The following
conclusions summarise our principal findings.

\vspace{0.5cm}

\noindent
\textbf{1. Exogenous features are the single largest improvement.}
Integrating weather (temperature, humidity, pressure, wind speed) and
calendar features reduces CRPS by approximately 20\%, directly
confirming the limitation identified in~\cite{smyl2025anyquantile} and
establishing exogenous integration as the highest-priority enhancement
for any future work on this benchmark. Importantly, the improvement
is robust: it persists across all multi-seed exogenous configurations
evaluated (CRPS range 168--174) and generalises across both Curriculum
v1 and v2 protocols. This single finding --- that meteorological inputs
matter more than architectural sophistication --- is the most impactful
practical contribution of this study.

\vspace{0.3cm}

\noindent
\textbf{2. Curriculum learning achieves the best multi-seed result.}
The two-stage Curriculum v2 protocol --- beta sampling to establish
tail calibration, followed by adaptive sampling to sharpen the
distribution --- achieves CRPS 172.21 ($\pm$1.5) with coverage 0.907.
This is the best reproducible result in the study and surpasses the
AQ-ESRNN benchmark (195.94) by approximately 12\%, despite the simpler
N-BEATS architectural backbone. The curriculum ordering is critical:
adaptive sampling applied without a calibrated initialisation
consistently degrades both CRPS and coverage.

\vspace{0.3cm}

\noindent
\textbf{3. CQR reliably corrects miscalibration under distribution shift.}
Applied post-hoc to a trained quantile model, CQR achieves 96.0\%
empirical coverage under a challenging winter-to-summer calibration
split, where the calibration and test distributions differ substantially.
The marginal CRPS cost (129.49 to 129.18) demonstrates that the
correction adds only minimal interval width, making CQR an efficient
and practical post-processing step for any deployed quantile forecaster.

\vspace{0.3cm}

\noindent
\textbf{4. Novel architectural heads establish future baselines.}
QCBC, TCR, and DBE all perform near the paper baseline (CRPS 212--217)
without exogenous features. DBE achieves a best-seed CRPS of 210.99,
marginally below baseline. While none of these contributions achieves
a clear improvement in isolation, they provide principled frameworks
and implementations for future research combining novel heads with
exogenous inputs.

\vspace{0.3cm}

\noindent
\textbf{5. Interaction and ordering effects are critical.}
The most dramatic negative result in the study --- monotonicity
constraints on an exogenous model, CRPS = 239.22 --- demonstrates
that enhancement strategies can have strongly negative interactions
when combined without curriculum ordering. This finding underscores
the value of the ablation-driven methodology adopted in this study.

\vspace{0.8cm}

\section*{Future Perspectives}

\noindent
Several directions emerge naturally from the findings of this work:

\vspace{0.3cm}

\noindent
\textbf{Ensemble of best configurations.}
Combining Stage 2-v2, the single-seed exogenous model (Exp 20), and a
series-embedding model via quantile-level median pooling could reduce
variance and push CRPS below any individual model. The diversity
of the three configurations (different training strategies, different
feature sets) suggests low inter-model correlation and thus strong
ensemble benefits.

\vspace{0.3cm}

\noindent
\textbf{Operational weather forecast inputs.}
The current study uses historical reanalysis data, which is a
retrospective reconstruction of atmospheric conditions. Operational
forecasting systems use \emph{weather forecasts}, which carry their
own uncertainty, particularly at the 24--48 hour lead times relevant
to day-ahead electricity markets. Replacing reanalysis inputs with
ensemble weather forecast products (e.g., ECMWF ENS or GFS) would
more realistically simulate the operational deployment scenario and
quantify the additional CRPS loss from weather forecast uncertainty.

\vspace{0.3cm}

\noindent
\textbf{Online CQR adaptation.}
The current CQR implementation uses a fixed calibration set and
therefore cannot adapt to concept drift (e.g., a gradual climate
shift or a structural change in demand behaviour). A sliding-window
calibrator that continuously updates the conformity score quantile
$\Delta$ using the most recent observations would maintain coverage
guarantees throughout the year without requiring a fixed historical
calibration period.

\vspace{0.3cm}

\noindent
\textbf{Transfer learning and per-country fine-tuning.}
The global cross-learning strategy trains a single model across all
35 countries simultaneously. An alternative is to pre-train globally
(for data efficiency and generalisation) and then fine-tune per country
on its own recent history. This approach could capture national demand
idiosyncrasies --- electricity market liberalisation, industrial
restructuring, residential electrification --- that the global model
can only partially learn through series embeddings.

\vspace{0.3cm}

\noindent
\textbf{Holiday and extreme event modelling.}
The current system treats public holidays as regular days. Prior work
has demonstrated that dedicated holiday models or rule-based
corrections significantly reduce residual error on atypical
days~\cite{son2022holiday, gunawan2021holiday, arora2018holiday}.
Integrating these corrections as a modular post-processing step,
alongside explicit extreme weather event flags (e.g., a temperature
threshold indicator for heatwave detection), would address the main
remaining source of tail miscalibration identified in this study.

\vspace{0.3cm}

\noindent
\textbf{Extension to other energy domains.}
The any-quantile framework and the curriculum learning methodology
developed in this study are architecture-agnostic and loss-agnostic.
Testing their transferability to related probabilistic forecasting
tasks --- solar photovoltaic generation, wind power output, natural
gas demand --- would demonstrate the generality of the contributions
and provide additional benchmarks for the community.
